% RESEARCH_DRAFT_SHA256: cbaaf3f12710c5ddf9bfa0412182bcd7cc8d66d60fa78aacf1f29853a94bfee0
% Layout and source-reference wrappers expanded for the web renderer.
\documentclass[11pt]{article}
\usepackage[margin=25mm]{geometry}
\usepackage{amsthm}
\usepackage{fontspec,unicode-math}
\setmainfont{Latin Modern Roman}
\setmathfont{Latin Modern Math}
\usepackage{amsmath,amssymb,mathtools,xurl,hyperref}
\hypersetup{colorlinks=true,urlcolor=blue,linkcolor=blue}
\setcounter{secnumdepth}{0}
\setlength{\parindent}{0pt}
\setlength{\parskip}{5pt}
\setlength{\emergencystretch}{3em}
\newcommand{\sref}[2]{\hyperlink{source:#1:#2}{[S#2]}}
\newcommand{\cataloguescope}[1]{\paragraph{Source record.}#1}
\newcommand{\floor}[1]{\left\lfloor#1\right\rfloor}
\newtheorem{proposition}{Proposition}
\newtheorem{lemma}[proposition]{Lemma}
\begin{document}
\paragraph{Research manuscript.} OpenAI \texttt{gpt-6-astra}, reasoning \texttt{ultra}; 11 October 2026.
\paragraph{Preserved source body.} The complete catalogue mathematical text and sources follow with layout and source links expanded, including source qualifications and its historical preparation notes. New research begins at the section headed Research attempt.

\section{Half-Integral Extreme Points for a Feedback-Vertex-Set Relaxation}\label{statement}

\subsection{Definitions and mathematical statement}

\paragraph{Shared notation.}Standard number systems are $\mathbb N,\mathbb Z,\mathbb Q,\mathbb R,\mathbb C$. Specialized notation and index ranges follow the cited source and the definitions below.


A feedback vertex set is a set of graph vertices whose removal leaves no cycles. A fractional relaxation replaces selected/unselected decisions by nonnegative numerical variables and valid linear constraints. An extreme point cannot be written as a proper convex combination of two other feasible points.
\paragraph{Record scope and qualifications.}
The catalogue title is retained as metadata. Named certificates, contracts, bounds or frameworks in that title are not asserted as results of the original source extract.
Source-record limitation: Duplicate half-integral extreme-point question for the same feedback-vertex-set relaxation retained under OWR ID 30004905.

Related catalogue records: \href{https://mehmetmars7.github.io/Erdosproblems-llm-hunter/problem.html?type=open_problems&id=30004905}{30004905}. Separate source IDs are preserved for their provenance.
\paragraph{Statement recorded in the supplied source.}
Does every basic feasible solution of the Chudak–Goemans–Hochbaum–Williamson linear-programming relaxation for a nontrivial feedback vertex set instance have a variable of value at least $1/2$? \href{https://doi.org/10.4171/owr/2024/50}{[S2]}
\paragraph{Source status.}
The supplied catalogue labels this source record open; that metadata was checked against the supplied snapshot on 2026-10-11. The label is preserved for provenance and is not an independent verification of present open status or mathematical correctness. The catalogue literature assessment is dated 2026-08-22.
\subsection{Short English statement}
The source asks about half-Integral Extreme Points for a Feedback-Vertex-Set Relaxation.
\subsection{Sources}
\begin{itemize}
\item[S1] Ulam AI and the UnsolvedMath Contributors, UnsolvedMath: A Curated Collection of Open Mathematics Problems, catalogue entry. Catalogue curation and metadata: CC BY 4.0. The source extract is retained; definitions, English summaries and qualifications are editorial additions.
\url{https://huggingface.co/datasets/ulamai/UnsolvedMath}.
\item[S2] Combinatorial Optimization (2025), Oberwolfach Reports, DOI 10.4171/owr/2024/50, pp. 2949--3000
\url{https://doi.org/10.4171/owr/2024/50}.
\item[S3] OWR report
\url{https://oa.tib.eu/renate/bitstreams/706cca58-da6b-4d8c-a979-ebf6e92658f0/download}.
\item[S4] FVS polyhedral results
\url{https://chekuri.cs.illinois.edu/papers/fvs-polyhedral-mpa.pdf}.
\item Source-specific Oberwolfach Reports licences, including CC BY-SA 4.0 where applicable, are retained; see THIRD\_PARTY\_NOTICES.md.
\end{itemize}
\paragraph{Preparation and review.}
The original source and UnsolvedMath contributors retain attribution for the source extract. Codex (GPT-6; exact serving revision not exposed) assisted with definitions, English summaries, scope notes and formatting. No new human mathematical authorship or independently verified solution is asserted. Catalogue identity checked against the complete supplied registry and current repository. Distinct upstream records, including repeated questions and broad or damaged extracts, are retained for complete coverage of the source CS category. A repeated record is not a new independent conjecture. Complete mathematical-equivalence and current-literature review has not been performed. Related canonical identities: 30004905. Preparation and final review follow contribution prompts C and D, with source gaps disclosed. Structural and rendering checks do not establish a new theorem or open status.
\label{end:statement}

\section{Research attempt}
\subsection{Literature, source context and scope}
The full catalogue body, including qualifications and historical source status, is preserved above. The original workshop report was inspected in the locally cached publisher PDF text. This attempt distinguishes the source question from any explicitly restricted model introduced below. Repeated records retain their own source IDs and do not become independent conjectures. Elementary results and attributed literature results are not claimed as new original-target progress.

\subsection{The actual relaxation and a symmetric candidate}
The primary paper specifies the strong-density relaxation
\[
 P_{\rm SD}(G)=\left\{x\ge0:
 \sum_{v\in S}(d_S(v)-1)x_v\ge |E[S]|-|S|+1
 \quad\text{for every }S\subseteq V\text{ with }E[S]\ne\varnothing\right\}.
\]
The question is the existence of a coordinate at least $1/2$ at every extreme point when the graph has a cycle. It does not say that every coordinate is half-integral. We tried the complete graph $K_n$, $n\ge3$, where symmetry suggests an all-small candidate.
\begin{proposition}
The vector $x_v=(n-1)/(2n)$ belongs to $P_{\rm SD}(K_n)$ and has every coordinate below $1/2$, but is not an extreme point.
\end{proposition}
\begin{proof}
For a set of size $s\ge2$, every degree is $s-1$ and the inequality is
\[
 (s-2)\sum_{v\in S}x_v\ge (s-1)(s-2)/2.
\]
For $s=2$ this is the zero inequality. For $s\ge3$ the displayed uniform point is feasible because $(n-1)/(2n)\ge(s-1)/(2s)$, with equality only when $s=n$.
To verify nonextremality explicitly, set $\varepsilon=1/(4n)$ and form $x^\pm=x\pm\varepsilon(e_1-e_2)$. They are distinct and positive, and their coordinates remain below $1/2$. The full-set sum is unchanged. For every proper subset with $s\ge3$, the original slack is $(s-2)(n-s)/(2n)\ge(s-2)/(2n)$, whereas the perturbation changes the left side by at most $(s-2)/(4n)$. Therefore both perturbed vectors are feasible and $x=(x^++x^-)/2$.
\end{proof}
The attempted symmetric counterexample therefore fails at precisely the extreme-point hypothesis. On $K_4$ the simple cycle-cover candidate $(1/3)^4$ also fails the full strong-density constraint, since $2\sum_vx_v=8/3<3$; replacing the named relaxation by cycle-cover inequalities would invalidate the attempt.
\subsection{Verification and remaining gaps}
All nonempty-edge subsets of $K_4$ were checked for the vector $(3/8)^4$ and its $\pm(1/16)(e_1-e_2)$ perturbations. The example explains why feasibility alone is insufficient. No proof controlling a full-rank family of tight strong-density constraints on general graphs has been found.

\section{References and source attribution}
The primary Oberwolfach workshop report, \url{https://doi.org/10.4171/owr/2024/50}, was inspected on 11 October 2026. Relevant source passages: 2024/50 pp. 2993--2994, Fiorini question. Catalogue and contributor attribution remain in the original Sources section. Mathematical arguments labeled as this attempt are proved above; cited external results remain attributed to their authors and are not independently reproved or claimed as our discoveries.

\url{https://chekuri.cs.illinois.edu/papers/fvs-polyhedral-mpa.pdf}. Chandrasekaran, Chekuri, Fiorini, Kulkarni and Weltge, Polyhedral aspects of feedback vertex set and pseudoforest deletion set. Opened primary 48-page paper; equation (1), the strong-density polyhedron, and Conjecture 1 on printed pages 3--4 were inspected. Its integrality-gap results are distinct from the unresolved extreme-point assertion.


\paragraph{Integration review.} Codex (GPT-6; exact serving revision not exposed) reviewed the written arguments and expanded layout and source-reference wrappers. Repository metadata and the production web renderer were checked using Python 3.13 and JavaScript with MathJax 3.2.2. This model-assisted review is not independent human proof verification; no PDF compilation or formal proof checking is asserted.

\section{Completion Estimate}
\noindent\textbf{COMPLETION: 0\%}
No substantiated new progress toward the entire original target is claimed. The worked results and failed route are delimited above; a known result, illustrative model, or auxiliary obstruction is not counted as original-target completion.

\section{Final status and attribution}
\textbf{Attempt status: unresolved.} This describes this research attempt, not a new ruling on the source catalogue's status or on external literature claims.
The mathematics and this self-audit were generated by OpenAI \texttt{gpt-6-astra}, reasoning \texttt{ultra}, on 11 October 2026. Codex (GPT-6; exact serving revision not exposed) orchestrated the batch and contributed formatting/repository review. No human mathematical authorship, independent human proof checking, or formal verification is claimed. The historical preparation paragraphs in the reproduced catalogue body concern that body only.

Contribution prompt A was applied to the complete definition and supplied source record, followed by the adversarial checks of prompt B. Local reproduction files are \texttt{batch011/checks.py}, \texttt{computations.json}, and the per-ID audit JSON. They were not public links at preparation time. Numerical checks are finite and do not prove unrestricted results. No LaTeX engine was found locally; compilation was not performed by this research agent.

\end{document}
